\documentclass[11pt]{article}
\usepackage[T1]{fontenc}
\usepackage{lmodern}
\usepackage[margin=0.9in]{geometry}
\usepackage{amsmath,amssymb,amsthm}
\usepackage{microtype}
\usepackage{xcolor}
\usepackage{booktabs}
\usepackage{needspace}
\usepackage{hyperref}
\hypersetup{colorlinks=true,linkcolor=black,citecolor=black,urlcolor=blue,
  pdftitle={Uniform bounds for greedy sequences avoiding consecutive sums},
  pdfauthor={Mark Vladimir Prizmic}}
\setlength{\emergencystretch}{2em}
\newtheorem{theorem}{Theorem}
\newtheorem{lemma}[theorem]{Lemma}
\newtheorem{corollary}[theorem]{Corollary}
\theoremstyle{remark}
\newtheorem{remark}[theorem]{Remark}
\newcommand{\N}{\mathbb{N}_{>0}}
\newcommand{\floor}[1]{\left\lfloor #1\right\rfloor}

\title{Uniform bounds for greedy sequences\\avoiding consecutive sums}
\author{Mark Vladimir Prizmic}
\date{November 27, 2026\quad\textnormal{\small Version 0.1}}

\begin{document}
\maketitle

\begin{abstract}
Fix an integer $r\ge2$. Starting with $a_1=1$, select each subsequent term
greedily, excluding every sum of $r$ consecutive previously selected terms.
We prove explicit bounds on both the selected sequence and its complementary
sequence of forbidden sums. With $C_r=r(r-1)/2-1$, the bounds are
\[
\floor{\frac{(r+1)n-C_r-1}{r}}
\le a_n\le
\floor{\frac{(r+1)n-1}{r}}\qquad(n\ge1).
\]
Consequently $a_n=(r+1)n/r+O_r(1)$, and the selected set has natural density
$r/(r+1)$. The case $r=4$ proves the upper bound recorded as Conjecture 1
in OEIS A394676 and supplies a lower bound. The proof is elementary and uses
the complementary partition together with strong induction. Exact-integer
verification code accompanies the note.
\end{abstract}

\section{The sequence and the result}

Let $r\ge2$ be fixed. Define $a_1=1$ and, once $a_1,\ldots,a_n$ have been
chosen, let $a_{n+1}$ be the smallest integer greater than $a_n$ outside
\begin{equation}\label{eq:forbidden}
\left\{\sum_{t=0}^{r-1}a_{j+t}:1\le j\le n-r+1\right\}.
\end{equation}
The set in \eqref{eq:forbidden} is empty when $n<r$. Thus $a_1,\ldots,a_r$
are $1,\ldots,r$. Throughout, consecutive terms means consecutive positions
in this sequence, not necessarily consecutive integers.

Hoang Xuan Thanh's OEIS entries A394576 and A394676 describe the cases
$r=3$ and $r=4$, respectively~\cite{oeis3,oeis4}. For $r=4$, the sequence
begins
\[
1,2,3,4,5,6,7,8,9,11,12,13,15,16,17,19,\ldots.
\]
For example, $10=1+2+3+4$ and $14=2+3+4+5$ are excluded.
Conjecture 1 in A394676 is the inequality
$a_n\le\floor{(5n-1)/4}$. We establish it as a special case of the following
general result.

\begin{theorem}\label{thm:main}
For an integer $r\ge2$, define $(a_n)$ by \eqref{eq:forbidden} and put
\[
b_j=\sum_{t=0}^{r-1}a_{j+t},\qquad C_r=\frac{r(r-1)}2-1.
\]
For every $j,n\ge1$,
\begin{equation}\label{eq:b-bounds}
(r+1)j\le b_j\le(r+1)j+C_r,
\end{equation}
and
\begin{equation}\label{eq:a-bounds}
\floor{\frac{(r+1)n-C_r-1}{r}}
\le a_n\le
\floor{\frac{(r+1)n-1}{r}}.
\end{equation}
\end{theorem}

\section{Complementarity and a floor identity}

\begin{lemma}\label{lem:partition}
The increasing sequences $(a_n)$ and $(b_j)$ partition the positive integers.
In particular,
\begin{equation}\label{eq:count}
a_i-i=\#\{j:b_j<a_i\}.
\end{equation}
\end{lemma}

\begin{proof}
At each finite stage, only finitely many integers are forbidden, so the greedy
construction continues indefinitely. The selected sequence is strictly
increasing and hence unbounded. Also,
\[
b_{j+1}-b_j=a_{j+r}-a_j\ge r>0.
\]
The sum $b_j$ exceeds its largest summand $a_{j+r-1}$ because $r\ge2$ and
all terms are positive. When that largest summand is selected, $b_j$ becomes
forbidden before the construction can reach it. Thus no forbidden sum is
selected. Conversely, every positive integer omitted by the greedy algorithm
was skipped because it was already in the finite forbidden set. This proves
the partition. Among $1,\ldots,a_i$, exactly $i$ integers are selected and
$a_i-i$ are forbidden, giving \eqref{eq:count}.
\end{proof}

\begin{lemma}\label{lem:floor}
For every integer $x$ and every positive integer $r$,
\[
\sum_{t=0}^{r-1}\floor{\frac{x+t}{r}}=x.
\]
\end{lemma}

\begin{proof}
Write $x=qr+s$ with $q\in\mathbb Z$ and $0\le s<r$. Of the $r$ terms in
the sum, $r-s$ equal $q$ and $s$ equal $q+1$. Their sum is $qr+s=x$.
\end{proof}

\section{Proof of the bounds}

\begin{proof}[Proof of Theorem~\ref{thm:main}]
Write $C=C_r$, so $C\ge0$. We first prove \eqref{eq:b-bounds} by strong
induction on $j$. The base case is
\[
b_1=1+\cdots+r=\frac{r(r+1)}2=(r+1)+C\ge r+1.
\]

Suppose $j\ge2$ and \eqref{eq:b-bounds} holds for every smaller index.
Fix $i\in\{j,\ldots,j+r-1\}$ and put $m=a_i-i$. By
Lemma~\ref{lem:partition}, $m$ counts the forbidden sums below $a_i$.
Since
\begin{equation}\label{eq:no-circularity}
b_j>a_{j+r-1}\ge a_i,
\end{equation}
we have $0\le m<j$. Consequently every forbidden value counted by $m$
has an index covered by the induction hypothesis.

\medskip\noindent\textit{Upper estimate for the selected terms.}
If $m\ge1$, then
\[
(r+1)m\le b_m<a_i.
\]
Substituting $m=a_i-i$ and using integrality gives
\[
r a_i\le(r+1)i-1,
\qquad
a_i\le U_i:=\floor{\frac{(r+1)i-1}{r}}.
\]
If $m=0$, then $a_i=i\le U_i$, so the same estimate holds.

\medskip\noindent\textit{Lower estimate for the selected terms.}
If $m+1<j$, the next forbidden value also has an index covered by the
induction hypothesis, and
\[
a_i<b_{m+1}\le(r+1)(m+1)+C.
\]
Again using integrality and substituting $m=a_i-i$ gives
\[
r a_i\ge(r+1)i-r-C.
\]
Equivalently,
\begin{equation}\label{eq:lower-i}
a_i\ge L_i:=\floor{\frac{(r+1)i-C-1}{r}}.
\end{equation}
Here we used $\lceil z/r\rceil=\floor{(z+r-1)/r}$ for integer $z$.

The remaining case is $m=j-1$. In this case $a_i=i+j-1$, and
\[
\begin{aligned}
r a_i-\bigl((r+1)i-r-C\bigr)
&=rj+C-i\\
&\ge(r-1)(j-1)+C\ge0,
\end{aligned}
\]
because $i\le j+r-1$. This proves \eqref{eq:lower-i} without appealing to
the as-yet unproved upper bound on $b_j$.

\medskip\noindent\textit{Summing the estimates.}
For integer $d$, Lemma~\ref{lem:floor} gives
\begin{equation}\label{eq:sum}
\begin{aligned}
\sum_{i=j}^{j+r-1}\floor{\frac{(r+1)i-d}{r}}
&=\sum_{i=j}^{j+r-1}\left(i+\floor{\frac{i-d}{r}}\right)\\
&=(r+1)j+\frac{r(r-1)}2-d.
\end{aligned}
\end{equation}
Taking $d=C+1=r(r-1)/2$ and $d=1$ shows that
\[
\sum_{i=j}^{j+r-1}L_i=(r+1)j,
\qquad
\sum_{i=j}^{j+r-1}U_i=(r+1)j+C.
\]
Since $L_i\le a_i\le U_i$, these identities establish
\eqref{eq:b-bounds} at index $j$ and complete the induction.

Finally fix arbitrary $n\ge1$ and put $m=a_n-n$. The now-global lower
bound on $b_m$ yields the upper bound in \eqref{eq:a-bounds} exactly as
above when $m\ge1$; the case $m=0$ is immediate. For the lower bound, the
next forbidden value $b_{m+1}$ always exists, and the now-global upper bound
gives $a_n<b_{m+1}\le(r+1)(m+1)+C$. The same rearrangement yields
$a_n\ge L_n$. This proves \eqref{eq:a-bounds} for every $n\ge1$.
\end{proof}

\section{Consequences and scope}

\begin{corollary}\label{cor:four}
For the sequence A394676, for every $n\ge1$,
\[
\boxed{\floor{\frac{5n-6}{4}}\le a_n\le\floor{\frac{5n-1}{4}}.}
\]
In particular,
\[
0\le\floor{\frac{5n-1}{4}}-a_n\le2,
\qquad
\lim_{n\to\infty}\frac{a_n}{n}=\frac54.
\]
\end{corollary}
\begin{proof}
Set $r=4$, for which $C_r=5$. The two floor bounds differ by at most
$\lceil5/4\rceil=2$. Dividing by $n$ gives the limit.
\end{proof}

\Needspace{8\baselineskip}
\begin{corollary}
For every fixed $r\ge2$, the selected set $A=\{a_n:n\ge1\}$ has natural
density $r/(r+1)$.
\end{corollary}
\begin{proof}
Equation~\eqref{eq:a-bounds} gives $a_n=\alpha n+O_r(1)$ with
$\alpha=(r+1)/r$. If $a_n\le x<a_{n+1}$, then
$\#(A\cap[1,x])=n$ and $x=\alpha n+O_r(1)$. Hence the ratio of this count
to $x$ tends to $1/\alpha=r/(r+1)$.
\end{proof}

For $r=2$, $C_r=0$ and the bounds coincide, giving
$a_n=\floor{(3n-1)/2}$. For $r=3$, they give
\[
\floor{\frac{4n-3}{3}}\le a_n\le\floor{\frac{4n-1}{3}}.
\]
This agrees with the three-term result recorded in Goyal's manuscript
collection~\cite{goyal}. The present argument treats all $r\ge2$ uniformly.
It does not determine the exact locations or lengths of the equality blocks
described in Conjecture 2 of A394676.

\section{Computation and research status}

The accompanying Python program generates terms directly from the greedy
rule using a queue of forbidden sums. Its generation step does not use the
bounds proved above. A separate verification pass checks the selected-term
bounds, the bounds on all completely generated forbidden sums, and the
complementary partition through the last selected term. Small cases are
also compared with an independent implementation of the definition.

For $r=4$, an exact run through $200\,000$ selected terms satisfies all
checks and gives $a_{200000}=249999$. The observed values of
$\floor{(5n-1)/4}-a_n$ are $0$, $1$, and $2$. The repository records the
complete run summary and a SHA-256 digest of the generated prefix so the
calculation can be reproduced. These finite checks support implementation
consistency; the argument in Section~3 proves the infinite statement.

\paragraph{Related work and priority.}
The source entries attribute the sequences to Hoang Xuan Thanh. Goyal's
repository records a proof for $r=3$ and identifies Thanh's related
Zenodo deposit as software~\cite{goyal,thanhsoftware}. The full contents of
that deposit were not available during preparation of this version. The
literature check is consequently incomplete. This note makes no claim of
priority or of being the first proof of the four-term or general result.
The status of an OEIS conjecture label alone does not establish novelty.

\paragraph{AI assistance and review.}
The proof, its generalization, the manuscript, and the verification code
were developed with substantial assistance from OpenAI's ChatGPT.
Additional AI-assisted reviews checked the induction and boundary cases.
This version has not received external human peer review and contains no
proof-assistant certification. The written proof is provided in full for
independent assessment.

\small
\begin{thebibliography}{9}
\bibitem{oeis4}
H.~X.~Thanh, \emph{A394676}, The On-Line Encyclopedia of Integer Sequences,
entry dated May 11, 2026. \url{https://oeis.org/A394676}.
Consulted November 27, 2026.

\bibitem{oeis3}
H.~X.~Thanh, \emph{A394576}, The On-Line Encyclopedia of Integer Sequences,
entry dated May 11, 2026. \url{https://oeis.org/A394576}.
Consulted November 27, 2026.

\bibitem{goyal}
D.~D.~Goyal, \emph{A greedy sequence that avoids sums of three consecutive
terms}, manuscript entry in the \texttt{oeis-proofs} collection, 2026.
\url{https://github.com/dvgyl/oeis-proofs/blob/main/CONTENTS.md}.
The manuscript map and repository metadata were consulted November 27, 2026.

\bibitem{thanhsoftware}
H.~X.~Thanh, \emph{Algorithm for OEIS A394576, A394676 and generalized
forbidden sums}, software deposit, Zenodo, 2026.
\url{https://doi.org/10.5281/zenodo.22674326}.
Bibliographic metadata identified in the related repository; full deposit
not inspected for this version.
\end{thebibliography}

\end{document}
